\section{Matchings}

\begin{definition}

    An edge set $M \subseteq E$ is called a \textbf{matching} in $G = (V, E)$ if no vertex is incident to
    more than one edge from $M$.
    \begin{itemize}
    \item A vertex $v$ is \textbf{covered} by $M$ if it is incident to an edge in $M$.
    \item A matching is \textbf{perfect} if every vertex is covered ($|M| = |V|/2$).
    \item A matching is \textbf{maximal} if it cannot be extended by adding edges.
    \item A matching is \textbf{maximum} if there is no matching with more edges.
\end{itemize}

\end{definition}

\begin{algorithm}

#Greedy Algorithm
#This leads to a maximal matching M. 
#The size of M is at least half the size of a maximum matching.

M = empty set
while E is not empty:
    choose an arbitrary edge e from E
    add e to M
    remove all edges incident to the endpoints of e from E
    
\end{algorithm}

\begin{theorem}

\textbf{Hall's Marriage Theorem}:\\

    A bipartite graph $G = (A \cup B, E)$ contains a matching of cardinality $|M|=|A|$ iff:
    $$\forall X \subseteq A : |X| \le |N(X)|$$
    
\end{theorem}

\begin{theorem}

\textbf{Frobenius Corollary}:\\

    Every $k$-regular bipartite graph contains a perfect matching.
    
\end{theorem}

\begin{definition}

    An \textbf{$M$-augmenting path} is a path that alternates between edges in $M$ and not in $M$,
    starting
    and ending in uncovered vertices.
    
\end{definition}

\begin{theorem}

\textbf{Berge's Theorem (1957)}:\\

    A matching $M$ is maximum iff there is no $M$-augmenting path.
    
\end{theorem}

